\section{Related Work}
\label{sec:related_conclusion}


\paragraph{Active Directory and PKI Security}
The BloodHound software models identity relationships for privilege-path investigation. SpecterOps~\cite{schroeder2021certified} systematized ADCS misconfigurations; Certipy~\cite{lyak2021certipy} supports certificate-service assessment. A restored configuration snapshot alone cannot reveal every intervening issuance. However, practical CA/KDC correlation is established: In.security describes SAN-abuse investigation in KQL~\cite{insecurity2023adcs}, and Splunk publishes an ESC1 authentication analytic~\cite{splunk2026esc1}. CIPG therefore does not claim novelty for joining issuance and authentication events. Its research question concerns certificate-specific identity consistency, possible versus definite answers, and revision under incomplete or conflicting evidence.

\paragraph{Graph learning and shortcut behavior}
Graph neural networks aggregate information over relationships. The original graph model and later reviews describe the assumptions and model families behind this approach~\cite{scarselli2008graph,wu2021survey,zhou2020review}. Convolutional graph models are reviewed by Wu et al.~\cite{wu2021survey}; the attention construction follows GAT~\cite{velickovic2018graph}. Relational graph convolution (R-GCN) and gated graph neural networks are alternative families reviewed in~\cite{wu2021survey,zhou2020review}; GATv2 changes the attention scoring order~\cite{brody2022attentive}. Our additive control is a matched ablation, not a comparison against all these families. Euler studies temporal link prediction for lateral movement~\cite{king2023euler}; its task differs from our resolved-permission predicate. Shortcut learning~\cite{geirhos2020shortcut} and adversarial robustness~\cite{wu2025understanding} motivate testing structural counterfactuals. Neither literature establishes that the bilinear gate used here improves classification. Our matched experiment finds no such improvement.

\paragraph{Provenance and incomplete evidence}
Surveys describe provenance capture, dependency growth and security analysis~\cite{zipperle2023survey,provenance2023security}. HADES addresses whole-network AD investigation~\cite{liu2026hades}, while CPD examines persistence~\cite{liu2026cpd}. These are related investigation systems, not executed baselines in this study. Broader provenance-IDS comparisons include HOLMES, KAIROS and FLASH; the recent evaluation-oriented synthesis~\cite{ipekbayrak2026evaluation} distinguishes their alerting and analysis units. These are literature comparisons, not reproduced baselines. Our identity-query experiment addresses a different task from end-to-end host intrusion detection; it supplies no ranking against those systems. Causality-preserving evaluation~\cite{suarez2025structure} further motivates keeping certificate-link correctness separate from attack reconstruction and detection. Established work on order-incomplete queries~\cite{amarilli2019computing}, temporal constraints~\cite{dechter1991temporal}, and record matching~\cite{elmagarmid2007duplicate} supplies broader reasoning foundations. Record matching alone does not guarantee correct attribution under compromised logs. Our contribution is the certificate-specific evidence contract and its evaluated boundaries, rather than a new general entity-resolution or temporal-query primitive. The timestamp-assumptions analysis~\cite{heinzl2026timestamps} is a preprint; it is not used as evidence of system performance.

\paragraph{Interdiction and remediation}
Attack-graph analysis includes reachability, cost and uncertainty models~\cite{zeng2019attackgraphs}. Goel et al. evaluate evolutionary policies for hardening AD graphs~\cite{goel2024hardening}. Classical network interdiction studies flow reduction and shortest-path disruption~\cite{wood1993interdiction,israeli2002shortest}. Security games provide a broader decision framework~\cite{tambe2011security}, while known cut problems provide the complexity foundation~\cite{garg1994multiway,dahlhaus1994multiterminal}. These objectives and sources are distinct. Our formulation adds assigned disruption costs and immutable edges to the evaluated instances. It does not establish a new general hardness result or certify an optimal remediation policy. The exhaustive small-instance comparison instead measures where the sequential heuristic falls short.

